\documentclass[12pt]{article}
\usepackage[utf8]{inputenc}
\usepackage[T1]{fontenc}
\usepackage{hyperref}
\title{Vera'a Language Resources for K-12 Curriculum Development}
\author{Compiled by Claude Sonnet 5 (claude-sonnet-5) for Andr\'as Kornai}
\date{}
\begin{document}
\maketitle
\section{Introduction}
Vera'a (ISO 639-3 vra; Glottocode vera1241) is an Oceanic (Austronesian) language of Vanua Lava in the Banks Islands, Vanuatu, with about 500 speakers. It is not to be confused with the neighbouring, closely related Vur\"es. The search found very little material beyond one grammar and a documentation project. No school materials were found.
\section{Primary Educational Resources}
None found in the search results. No primers, textbooks or school programmes for Vera'a were found.
\section{Linguistic and Reference Materials}
\begin{itemize}
\item Stefan Schnell, \emph{A grammar of Vera'a} (2011, Ph.D. dissertation, Kiel University). Source record: \url{https://tlopo.clld.org/sources/Schnell2011}
\item Wikipedia overview of the language: \url{https://en.wikipedia.org/wiki/Vera%CA%BCa_language}
\item Grambank entry: \url{https://grambank.clld.org/languages/vera1241}
\end{itemize}
No Vera'a dictionary was found. The Vur\"es dictionary and grammar by Catriona Malau concern a different language and are not Vera'a resources.
\section{Religious and Community Resources}
None found in the search results.
\section{Digital and Audio Resources}
A DOBES documentation project (2006--2012) covered Vera'a and Vur\"es; a poster describing it is at \url{https://dobes.mpi.nl/posters/print/oceania-vures-print.pdf}. The search results did not confirm what recordings are publicly accessible.
\section{Curriculum Development Context}
Vera'a is vigorous as a first language, but the search found no evidence of school use. A curriculum would have to start from the Schnell grammar and the DOBES recordings, and would need community-led orthography and materials development.
\end{document}